\section{Accessing Models with \texttt{env[...]}}

\begin{frame}[fragile]{The Pattern \texttt{self.env['model.name']}}
	\begin{itemize}
		\item This is the standard way to get a model from the environment.
\begin{block}{Method Syntax}
\begin{lstlisting}
Model = self.env['student.student']
records = Model.search([])
\end{lstlisting}
\end{block}
		\item \texttt{self.env['student.student']} returns an empty recordset of that model in the current environment.
		\item From there you call ORM methods: \texttt{create}, \texttt{search}, \texttt{browse}, \ldots
	\end{itemize}
\end{frame}

\begin{frame}[fragile]{Practical Examples}
\begin{lstlisting}
# Search students
students = self.env['student.student'].search([('active', '=', True)])

# Create a classroom
classroom = self.env['school.classroom'].create({
	'name': 'Grade 12-A',
})

# Browse by ID
partner = self.env['res.partner'].browse(7)
\end{lstlisting}
	\begin{itemize}
		\item Always use the technical model name (\texttt{\_name}), not the Python class name.
	\end{itemize}
\end{frame}

\begin{frame}[fragile]{Same Model, Different Environments}
\begin{lstlisting}
students = self.env['student.student']
admin_students = self.env['student.student'].sudo()
company_students = self.env['student.student'].with_company(company)
\end{lstlisting}
	\begin{itemize}
		\item All three access the same model.
		\item But rights, company, or context may differ.
		\item This is why environment control is so important.
	\end{itemize}
\end{frame}
